\documentclass[11pt]{article}
\usepackage[margin=1in]{geometry}
\usepackage{titlesec}
\usepackage{enumitem}
\usepackage{xcolor}
\usepackage{fancyhdr}
\usepackage{tcolorbox}
\tcbuselibrary{skins}

\definecolor{isaziteal}{HTML}{2AB6C4}
\definecolor{isazipink}{HTML}{C7327A}
\definecolor{isaziorange}{HTML}{F0803C}
\definecolor{isazigold}{HTML}{F0B429}

\titleformat{\section}{\Large\bfseries\color{isazipink}}{Slide \thesection}{0.6em}{}
\titleformat{\subsection}{\normalsize\bfseries\color{isaziteal}}{}{0em}{}

\newtcolorbox{visual}{colback=isaziteal!8,colframe=isaziteal,fonttitle=\bfseries,title=VISUAL / CANVA DIRECTION,boxrule=0.8pt,arc=2pt}
\newtcolorbox{scriptbox}{colback=gray!5,colframe=gray!40,fonttitle=\bfseries,title=WHAT I SAY (script),boxrule=0.6pt,arc=2pt}
\newtcolorbox{onslide}{colback=isazigold!10,colframe=isaziorange,fonttitle=\bfseries,title=ON-SLIDE TEXT (keep it this short),boxrule=0.8pt,arc=2pt}

\pagestyle{fancy}
\fancyhf{}
\fancyfoot[C]{\small\color{gray}Sibusiso Khumalo -- Isazi Vac Program Presentation}
\fancyfoot[R]{\small\color{gray}\thepage}
\renewcommand{\headrulewidth}{0pt}

\title{\color{isazipink}\Huge Isazi Vac Program\\\Large Presentation Script}
\author{Sibusiso Khumalo}
\date{}

\begin{document}
\maketitle

\begin{tcolorbox}[colback=isazipink!6,colframe=isazipink,title=HOW TO USE THIS DOCUMENT,fonttitle=\bfseries]
This is your talking script, organized to match Isazi's 7-slide vac-program template exactly.
For each slide you get three things: the \textbf{on-slide text} (short -- this is all that
should appear in Canva), the \textbf{visual direction} (what to actually design), and
\textbf{what I say} (your spoken script, delivered from memory, not read off the slide).
Rule of thumb: if a sentence is in the "what I say" box, it does NOT belong on the slide.
\end{tcolorbox}

% ---------------------------------------------------------------
\section{Identity}

\begin{onslide}
Sibusiso Khumalo\\
Clinical Medicine $\rightarrow$ Electrical Engineering $\rightarrow$ Isazi.ai\\
\textit{Still exploring where I fit in tech}
\end{onslide}

\begin{visual}
Full-bleed photo of you, left third of the slide. Right two-thirds: a simple horizontal
timeline/arrow graphic with three icons -- a stethoscope, a circuit/resistor symbol, and
the Isazi logo dot-mark -- connected by an arrow flowing left to right. Use Isazi's
gradient (teal $\to$ pink $\to$ orange $\to$ gold) as the arrow's color, echoing the logo.
Keep text minimal; the timeline graphic tells the story.
\end{visual}

\begin{scriptbox}
Hi everyone, I'm Sibusiso Khumalo. My path here has not been a straight line, and I think
that's actually the point. I started out studying clinical medical practice -- I wanted
to help people directly, one on one. What that degree really gave me was something I
didn't expect: a deep education in reading people, listening carefully, and communicating
under pressure. Those are human-interaction skills, not technical ones, but they turn out
to matter enormously in engineering too. I then made the switch to electrical engineering,
because I realized what I actually loved wasn't medicine specifically -- it was solving
hard, structured problems. I landed at Isazi almost by accident, through a hackathon,
and I stayed because of the culture: this is a place that incubates problem-solvers,
not just people who already have all the answers. I'm still figuring out exactly where
I fit in tech long-term, and that's what this month was about.
\end{scriptbox}

% ---------------------------------------------------------------
\section{Why Electrical Engineering, and Why Isazi}

\begin{onslide}
From bedside skills to circuit boards\\
Same instinct: understand the system, find what's broken, fix it
\end{onslide}

\begin{visual}
Split-panel slide, two halves separated by a thin gradient line. Left half: a soft
illustration/icon of a doctor's stethoscope with a small "listening / empathy" motif.
Right half: a circuit-board or PCB icon. In the middle, overlapping, a single lightbulb
or magnifying-glass icon to represent the shared thread: problem-solving. Very little
text -- let the icon pairing do the talking.
\end{visual}

\begin{scriptbox}
What drew me to electrical engineering was realizing that medicine had taught me how to
diagnose problems in people, and engineering let me diagnose problems in systems --
same muscle, different subject matter. I wanted to test that in the real world, not just
in a classroom, and Isazi was the obvious place to do that: it's a company that builds
real AI products for real businesses, under real constraints, not toy problems. I wanted
one month where the thing I built either worked in production or it didn't -- no
in-between, no simulated grading. That's what pulled me here specifically, over a
more conventional internship.
\end{scriptbox}

% ---------------------------------------------------------------
\section{Problem Statement}

\begin{onslide}
Well-connected people are sitting on gold mines of contacts\\
they don't have time to open.
\end{onslide}

\begin{visual}
A single bold stat as the hero visual: something like ``50{,}000+ contacts. 0 minutes to
spare.'' rendered huge, center-screen, in the Isazi gradient. Underneath, a small
supporting icon row: a clock, a WhatsApp bubble, and a tangled-network icon. This slide
should feel like a punch, not a paragraph -- one number, one line, done.
\end{visual}

\begin{scriptbox}
Here's the problem. Think about any well-connected professional -- a senior banker, an
investor, a partner at a firm. Over a career they build up a network of thousands of
real, high-value contacts. Constantly, friends and colleagues text them: ``hey, do you
know a good corporate lawyer in London?'', ``can you connect me with someone in private
equity?''. Each one of those messages is small on its own. But multiply it by dozens of
requests a week, and the real cost isn't goodwill -- it's time. Every request means
stopping what you're doing, mentally scanning your network, remembering who's relevant,
writing a warm, personal introduction, and doing it in a way that doesn't damage the
relationship on either side. For someone whose time is worth the most, that's exactly
the task they can least afford to do manually, and exactly the task nobody has properly
automated -- because it needs judgment, not just a search bar. That's the gap I built
into: give that person back their time, without losing the personal touch that makes
an introduction actually work.
\end{scriptbox}

% ---------------------------------------------------------------
\section{Project -- The Elevator Pitch}

\begin{onslide}
AI Middleman: your WhatsApp, minus the busywork.\\
\textit{[Live demo follows]}
\end{onslide}

\begin{visual}
Three-step horizontal flow diagram, icon-driven: (1) WhatsApp message bubble in $\to$
(2) a small "brain / gears" icon labeled AI $\to$ (3) an approve/checkmark icon out to a
second WhatsApp bubble. Keep the whole diagram to icons plus 2--3 word labels
(``Friend asks'', ``AI matches \& drafts'', ``You approve''). No paragraphs. This is the
slide where you switch to the live demo, so leave room at the bottom for a
``Watch this live $\to$'' cue.
\end{visual}

\begin{scriptbox}
In thirty seconds: AI Middleman is a system that sits inside WhatsApp and does the
matching and drafting work for you. A friend texts a request in plain language -- in
English or Afrikaans, typos and all. The system figures out whether it's actually a
contact request, works out who in the network is the best fit using a two-stage
AI pipeline, writes a warm introduction in your own voice, and only then hands it to you
to approve, edit, or skip with one tap. You stay fully in control of every introduction
that goes out -- the AI never sends anything without you. Let me just show you, live.
\textbf{[switch to live demo here]}
\end{scriptbox}

% ---------------------------------------------------------------
\section{What I Learned}

\begin{onslide}
Technical: production AI is a pipeline, not a prompt.\\
Human: the fastest way to unblock yourself is to ask.
\end{onslide}

\begin{visual}
Two stacked cards/icons side by side, labeled ``Technical'' and ``Human''. Technical card:
a small pipeline/funnel icon (representing the keyword-filter $\to$ LLM-match $\to$
draft stages). Human card: a speech-bubble/question-mark icon. Keep each card to a
single short phrase; let the icons carry the visual weight.
\end{visual}

\begin{scriptbox}
Technically, the biggest lesson was that a single clever prompt is never enough for a
real product -- you need a pipeline. I ended up with a two-stage matching system: a
cheap, deterministic keyword filter that narrows tens of thousands of contacts down to a
shortlist, and only then a more expensive LLM call to actually rank and reason about
that shortlist. That split is what makes the system fast and cheap enough to run on
every message, instead of an expensive model call trying to search the entire dataset
every time. On the human side, the real takeaway was about asking for help early. It's
tempting to sit with a bug for hours because you feel like you should be able to solve
it yourself -- but the fastest debugging tool here was always another person's fresh
eyes, and the culture at Isazi actively rewards asking rather than quietly struggling.
That's a mindset shift, not a technical one, but it's the one I'll keep using longest.
\end{scriptbox}

% ---------------------------------------------------------------
\section{What Got Stuck}

\begin{onslide}
The AI kept mixing up isiZulu and Afrikaans.\\
Honest fix: a hard rule beats a hopeful prompt.
\end{onslide}

\begin{visual}
A ``before / after'' two-column visual. Left column (red-tinted): a WhatsApp bubble
showing a garbled/wrong-language reply, small "X" icon. Right column (green-tinted):
the same request with a correct reply, small checkmark icon. This is your most authentic,
memorable slide -- let the visual literally show the bug and the fix, no abstract icons.
\end{visual}

\begin{scriptbox}
I want to be honest about something that genuinely broke. I'd restricted the system's
language detection to English and Afrikaans, because a broader multi-language attempt
was unreliable and I didn't want to embarrass myself or the friend messaging in
during the actual demo. But that fix introduced a new bug: because Afrikaans became the
only non-English option, the AI started force-fitting real isiZulu messages into
``Afrikaans'' too, since to a small model, both just look like ``not English.'' I only
caught this by actually reading through live conversation logs, not by assuming things
were fine. The fix wasn't a smarter prompt -- it was a deterministic rule: a short list
of words and prefixes that only ever appear in isiZulu, checked in code, that overrides
the AI's guess whenever it's about to get it wrong. The lesson I'm taking from this is
that when correctness really matters, you don't ask the model to be careful -- you build
a hard backstop around it.
\end{scriptbox}

% ---------------------------------------------------------------
\section{Thank You}

\begin{onslide}
Thank you.
\end{onslide}

\begin{visual}
Minimal, warm closing slide. Isazi logo centered, small, on a clean background (use the
dark logo variant if the room lights are dimmed, light variant otherwise). Just the
words ``Thank you'' beneath it, generous white space. No bullet points, no clutter --
this slide should feel like a breath, not more information.
\end{visual}

\begin{scriptbox}
Thank you -- genuinely. This month gave me something I didn't expect: proof that I can
take an idea, ship it, break it in front of real people, and fix it, all in a few weeks.
\textit{[Personalize by name here: thank the specific mentor(s) who reviewed your code
or unblocked you, the person who answered your questions when you were stuck on the
language bug or the matching pipeline, and anyone who made you feel like asking
"dumb" questions was safe. Say their names. Be specific about one moment they helped --
that specificity is what makes a thank-you land instead of sounding generic.]}
I'm walking away from this program with a real product, a lot of new judgment about
what "production-ready" actually means, and a much clearer sense of where I want to
grow next. Thank you for the month, and thank you for the space to figure this out.
\end{scriptbox}

\end{document}
